\section{Forcing an entire completed divisor fiber into its ordinary product}
\label{sec:r13-divisor-fiber-product-localization}

Use the algebraic standard etale divisor chart of DT and the
fixed surface preparation of UA.  This section adds finitely
many central coordinate functions.  Their purpose is to
recognize the selected product by inequalities on the entire
moment fiber.  No assumption on the moment rank of another
point in that fiber is made.

Write \(T=(T_1,T_2)\), let \(\Theta\) be the selected etale
coordinate on the original threefold chart, and let
\(\Theta_0\) be the surface coordinate.  The fixed equations
are
\[
 Q(Z,T,\Theta)=0,\qquad P(T,\Theta_0)=0,
 \qquad Q(0,T,\Theta)=P(T,\Theta)R_0(T,\Theta).
\tag{UF1}
\]
All coefficients lie in \(R[T,G^{-1}]\), after including
finitely many selected localizers.  The functions \(T_i\),
\(\Theta_0\), the relevant localizers, and the selected
derivative and complementary factors have norm one at the
surface point \(v\).  This is the zero-slice chart used in
DT.  Passing to a smaller such chart or multiplying a
coordinate by an original constant does not change the
construction.  In a degree-one or trivial etale presentation,
if the selected \(\Theta_0\) has zero residue, replace both
\(\Theta_0\) and \(\Theta\) by their sum with one.  This
integral coordinate change makes their core norms one and
preserves UF1 and its derivative conditions.

\subsection{Rational coordinates which recognize a surface neighborhood}

\begin{lemma}[A finite rational neighborhood packet]
\label{lem:r13-rational-neighborhood-packet}
Let \(W\) be an ordinary neighborhood of \(v\) in this
surface chart.  Let \(h_1,\ldots,h_s\) be finitely many
actual surface units of norm one at \(v\).  There are a
smaller ordinary neighborhood \(W'\subset W\), rational
functions \(r_j\) with
\[
                         |h_j|=|r_j|\quad\hbox{on }W',
\tag{UF2}
\]
and finitely many pairs of integral polynomial functions
\(A_a(T,\Theta_0),B_a(T,\Theta_0)\), all of norm one
at \(v\), such that any surface point at which all the
selected chart localizers have norm one and
\[
                       |A_a|=|B_a|=1\quad\hbox{for every }a
\tag{UF3}
\]
belongs to \(W'\) and satisfies \(|h_j|=1\) for every
\(j\).  Denominators involving the fixed localizers are
cleared in the displayed polynomials.
\end{lemma}

\begin{proof}
SA, or its single-function residue approximation step,
replaces each \(h_j\) by a rational surface function
with the same norm on one ordinary neighborhood of
\(v\).  Intersect these finitely many neighborhoods
with \(W\).  The topology of an affine algebraic
analytification is generated by norms of its regular
functions.  Choose a finite basic algebraic open
neighborhood inside that intersection.  Include upper
bounds for a finite affine coordinate-generating tuple
if a fixed bounded chart is desired.  Every nonzero
regular function has nonzero value at the generic
model point \(v\); its value is in \(|K^\times|\).
Normalize it by an original constant to have norm
one.  Retaining its norm-one equality implies all
the strict inequalities defining that basic open.

Include these finitely many normalized rational
functions and the \(r_j\).  Each has norm one at
the generic special-fiber point.  It can be written
as a quotient of two functions in a fixed smooth
integral affine model ring with nonzero reductions.
Indeed cancel the common \(\pi\)-order of its
numerator and denominator in the localization at
that generic point, and add finitely many
nonzero-reduction localizers.  Clearing these
localizers gives integral polynomial numerator
and denominator representatives, each with norm
one at \(v\).  Retain both.  Their norm-one
equalities give the required rational-function
equalities, placing the point in the chosen basic
open and giving \(|r_j|=1\).  Call that open
\(W'\).  Equation UF2 proves the last assertion.
\end{proof}

For each polynomial in this packet, define its direct
central lift by substituting \(\Theta\) for
\(\Theta_0\), without changing its coefficients:
\[
       \widehat A_a=A_a(T,\Theta),\qquad
       \widehat B_a=B_a(T,\Theta).
\tag{UF4}
\]
Their residues are nonzero along \(D\), so these
are norm-one actual units on a common ordinary
neighborhood of the original central point.

\begin{lemma}[The entire fiber has a selected zero-slice point]
\label{lem:r13-entire-fiber-zero-slice}
There is a finite actual central tuple \(J_D\)
containing \(Z\), the chart coordinate and localizer
functions, the direct lifts in UF4, and the selected
polynomial derivative, with the following property.
On a sufficiently outer collar, if an original point
\(x\) in their common ordinary chart satisfies
\[
                    t_{J_D}(x)=t_{J_D}(q_D(r)),
                             \qquad r<1,
\tag{UF5}
\]
then it belongs to the selected ordinary product
tube and its zero-slice point \(x_0\) belongs to
\(W'\) and satisfies \(|h_j(x_0)|=1\) for all \(j\).
This assertion includes arbitrary moment ranks and
all commensurable radial values.
\end{lemma}

\begin{proof}
If \(P(T,\Theta)\) is not identically zero on the
central chart, retain it as another actual unit,
after original-constant normalization if necessary.
It is nonzero at the central generic point, so
its unit domain contains one ordinary central
neighborhood.  If it is identically zero, no
coordinate is needed for it.  Retain
\(P_\Theta(T,\Theta)\), and the finite zero-slice
and chart localizers, in addition to UF4.  Every
norm-one function in this list is indeed norm
one at \(q_D(r)\): on the product,
\(|\Theta-\Theta_0|<1\), and integral polynomial
evaluation preserves every nonzero reduction.
Also
\[
 |P(T,\Theta)(q_D(r))|<1,
 \qquad |P_\Theta(T,\Theta)(q_D(r))|=1.
\tag{UF6}
\]
This is the selected DT Hensel comparison.  The
finite collar merely places all points and unit
domains in the original central chart.

Equation UF5 gives exactly these inequalities
at \(x\), together with \(|Z(x)|=r<1\),
\(|T_i(x)|=|\Theta(x)|=1\), and norm one for
every retained localizer.  Coefficients of \(P\)
therefore have norm at most one.  Hensel's lemma
in the complete field \(\mathcal H(x)\) gives
the unique root \(\theta_0\) with
\[
                       |\theta_0-\Theta(x)|<1.
\tag{UF7}
\]
If \(P(T,\Theta)\) is identically zero, take
\(\theta_0=\Theta(x)\).  This root defines
the surface point \(x_0\); all selected
localizers remain nonzero there.

Integral polynomial evaluation in UF7 gives
\[
 |A_a(T,\theta_0)-\widehat A_a(x)|<1,
 \qquad
 |B_a(T,\theta_0)-\widehat B_a(x)|<1.
\tag{UF8}
\]
Their central norms at \(q_D(r)\) are one,
so UF5 and UF8 give UF3 at \(x_0\).  The
preceding lemma gives \(x_0\in W'\) and
\(|h_j(x_0)|=1\).

Finally the polynomial \(Q\) has a unique
selected root within distance less than one
of \(\theta_0\).  Its Newton inequalities
follow from UF1, \(|Z|<1\), the retained
derivative and complementary factor norms,
and the same integral coefficient bounds.
The original \(\Theta(x)\) is that root.
The values \(Z(x),T(x),\theta_0\) in
\(\mathcal H(x)\) specify a joint seminorm,
and hence an \(\mathcal H(x)\)-valued
point of the analytic product.  Its two
underlying factor points alone would not
specify that product point.  Locally these
root constructions are analytic: continuity
gives strict Newton margins, and iteration
converges on an affinoid neighborhood.  Root
uniqueness glues the resulting inverse
analytic map to the product.  Composing
with the selected product map fixes
\(Z,T\), and fixes \(\Theta\) by the
preceding unique-root argument, so the
composition is the identity.  Thus
the conclusion is membership in an actual
ordinary product domain.  No graded-rank
count has been used.
\end{proof}

\subsection{The complete central image near the supplied ray}

Take \(H\) to be the fixed final protected surface
tuple in the proof of UA.  Its unit domain is fixed
before the annular interval.  For each interval
\(I\), the NSA construction is valid on the full
inverse image of a target neighborhood of
\(t_H(v)\) in that domain.  In particular it
contains the entire central \(H\)-fiber in any
fixed smaller unit domain.  This is the final
GP complete-image conclusion of NSA, retained
in UA's fixed geometric preparation.

Choose \(W\) inside that fixed unit domain,
and apply the packet lemma to every coordinate
of \(H\).  Include direct central representatives
of the rational functions \(r_j\) as well as
their separate numerators and denominators in
\(J_D\).  Near every point of the fiber in
UF5, their norms equal those of the surface
\(H\)-coordinates, by UF2 and the strict
polynomial estimates UF8.

\begin{theorem}[Completed core-ray fiber localization]
\label{thm:r13-completed-ray-fiber-localization}
For the fixed divisor and fixed central finite-PS
tensor in UA, there is one finite central
refinement \(F\), and a protected ordinary
central domain, such that for every compact
interval \(I\) in one sufficiently outer
collar the following hold.
\begin{enumerate}
\item The entire \(F\)-moment fiber over every
\(t_F(q_D(r))\), \(r\in I\), is contained
in an ordinary domain carrying the transported
NSA covector.  The same assertion holds for
every actual source in these fibers, including
points whose earlier moment rank is only two
or less.
\item There is an open target neighborhood
\(\mathcal O_I\) of
\(\{t_F(q_D(r)):r\in I\}\) whose complete
ordinary inverse image lies in that domain.
On its complete image the transported finite-PS
covector has the full all-face strain identity.
It has the fixed-control approximation and
all-background conclusions of UA for this
interval.
\end{enumerate}
One finite refinement works for a supplied
finite list of divisors and their respective
core rays.  The statement concerns neighborhoods
of those completed core rays; it does not
assert that they include every new ray which
projects to an earlier ray.
\end{theorem}

\begin{proof}
Retain \(J_D\), all original coefficient
coordinates, and the finite rational
representatives just specified.  Apply one
central GP on their common original ordinary
domain.  Call the completed tuple \(F\)
and its compact protected closure \(D_0\).
The protecting coordinates are constant on
an ordinary neighborhood of the original
central point, so every sufficiently short
core ray is retained unchanged.  All core
values in a sufficiently outer collar lie
inside the open protected target box.

Let \(x\) be any point in the entire
\(F\)-fiber over a core value.  Projection
to \(J_D\) gives UF5.  The preceding lemma
places its zero-slice point in the fixed
surface unit domain and in the entire
central \(H\)-fiber.  UA's fixed-control
lemma puts that full protected fiber inside
the valid AF product for the chosen radial
interval.  The final NSA primitive is valid
there on the complete image.  The direct
central lifts and UF8 transfer it on an
ordinary neighborhood of \(x\).  The
surface covector is a single ambient
formula \(\Lambda_I\); these transferred
formulas therefore agree wherever their
domains overlap.  Take their union as the
ordinary domain in the first assertion.
This argument has not removed a source
according to its rank.

The inverse image in \(D_0\) of the
compact set of core values for \(I\)
is compact and, by the first assertion,
is contained in that ordinary domain.
The closed complement in \(D_0\) has
compact moment image, disjoint from those
core values.  An open target neighborhood
avoiding this compact image and contained
in the original protected box is the
required \(\mathcal O_I\).  Protection
ensures that its full inverse image is
computed inside \(D_0\), and hence lies
in the stated ordinary domain.

The coefficient identity was transported
on every face, so it holds on this entire
inverse image and on every later retaining
refinement, including rank degeneration.
Apply the same strain smoothing argument
as in UA on a smaller protected truncation.
For finitely many supplied divisors, take
the union of the finite packets before
the final central GP and the maximum of
their finitely many radial thresholds.
Projection to each packet proves its
fiber assertion separately.  No further
iteration of surface preparation is
required for this finite list.
\end{proof}

